% ID: SEL-D0016-BD-E07
% ESTADO: validada. S017. Revisión matemática, didáctica y visual.
\textbf{a) Centro y semiamplitud.}\par\nopagebreak[4]
El intervalo es simétrico respecto de la media muestral. Su centro y su semiamplitud son
\[
\boxed{\overline x=\frac{153.65+162.35}{2}=158\text{ km}},\qquad
\boxed{E=\frac{162.35-153.65}{2}=4.35\text{ km}}.
\]

\textbf{b) Efecto de aumentar $n$.}\par\nopagebreak[4]
La desviación poblacional es $\sigma=\sqrt{225}=15$ km. Con confianza fija, el cuantil $z$ se mantiene, y $E=z\sigma/\sqrt n$ disminuye al aumentar $n$.

\textbf{c) Confianza del 95\,\% y error inferior a 3 km.}\par\nopagebreak[4]
Usamos $z_{0.975}\approx1.959963985$. La desigualdad estricta es
\[
z_{0.975}\frac{15}{\sqrt n}<3
\quad\Longleftrightarrow\quad n>25z_{0.975}^2\approx96.036470517.
\]
Por tanto, $\boxed{n_{\min}=97}$. El error con $96$ supera $3$ km y con $97$ queda por debajo. Usar $1.96$ da el mismo entero.
